% ============================================================
% AstraHeal Paper 2
% Evidential Uncertainty-Aware Fault Diagnosis for Autonomous
% Spacecraft Health Management
% Rewritten LaTeX source based on the supplied manuscript.
% ============================================================

\documentclass[conference]{IEEEtran}

\usepackage{amsmath,amssymb}
\usepackage{booktabs}
\usepackage{multirow}
\usepackage{array}
\usepackage{siunitx}
\usepackage{url}
\usepackage{graphicx}
\usepackage{cite}

\title{Evidential Uncertainty-Aware Fault Diagnosis for Autonomous Spacecraft Health Management}

\author{
\IEEEauthorblockN{AstraHeal Research Group}
\IEEEauthorblockA{
Autonomous Systems \& Aerospace Research\\
AstraHeal Paper 2 --- Formal Scientific Research Series\\
Repository: \url{https://github.com/madankalyan2211/AstraHeal}
}
}

\begin{document}

\maketitle

\begin{abstract}
Autonomous spacecraft must continue to make health-management decisions when
communication with ground operators is delayed or temporarily unavailable.
This creates a particular problem for fault diagnosis: a system must not only
identify a known failure, but also recognize when the available telemetry does
not provide enough evidence for a reliable diagnosis. Conventional rule-based
Fault Detection, Isolation, and Recovery (FDIR) systems are generally
conservative, while conventional machine-learning classifiers can remain
highly confident when presented with failures that were not represented during
training.

This work develops an evidential diagnostic architecture for spacecraft health
management. The proposed method represents competing fault hypotheses with a
Dirichlet distribution and uses Mahalanobis distance on physically motivated
telemetry features to determine the amount of evidence assigned to each
hypothesis. Epistemic uncertainty is used to represent lack of knowledge about
an observation, whereas aleatoric uncertainty captures ambiguity associated
with noisy telemetry.

The evaluation contains five controlled studies covering known-fault
classification, uncertainty behavior, out-of-distribution (OOD) detection,
telemetry-noise robustness, and component ablation. The benchmark includes
1,200 known-fault evaluations, 1,400 uncertainty-regime tests, 1,200 OOD
frames, 2,100 telemetry-noise sweeps, and systematic component removals.
On held-out known-fault frames, the proposed system obtains a Macro-F1 of
0.9533 with an Expected Calibration Error (ECE) of 0.0094. Under increasing
telemetry noise, aleatoric uncertainty increases while epistemic uncertainty
remains bounded for in-distribution observations. For previously unmodeled
faults, epistemic uncertainty increases substantially, producing an OOD
AUROC of 0.9422 and AUPRC of 0.9516. Maximum Softmax Probability (MSP),
in contrast, obtains an AUROC of 0.4313. The experiments also expose two
important limitations: cancellation in the metric representation during
compound failures and reduced sensitivity to sensor-polarity reversals when
magnitude-only features are used.
\end{abstract}

\begin{IEEEkeywords}
autonomous spacecraft, fault diagnosis, evidential deep learning,
Dirichlet distribution, epistemic uncertainty, aleatoric uncertainty,
out-of-distribution detection
\end{IEEEkeywords}

\section{Introduction}

Space missions increasingly require onboard autonomy because a spacecraft
cannot always depend on an immediate ground response. In Low Earth Orbit
(LEO), orbital geometry can interrupt line-of-sight communication for extended
periods, while deep-space missions introduce communication delays ranging from
minutes to hours. During such intervals, events involving the electrical power
system, thermal subsystem, or battery can develop faster than a ground
operator can diagnose and respond to them.

Traditional Fault Detection, Isolation, and Recovery (FDIR) architectures
address this problem with carefully designed thresholds, rules, and safe-state
transitions \cite{chien2005,muscettola1998,williams1996}. Such mechanisms are
valuable because their behavior is predictable. However, an overly conservative
response can also force the spacecraft into Safe Mode for situations that do
not actually threaten the vehicle. Safe Mode can interrupt payload operations,
discard observation opportunities, and reduce mission productivity.

Machine-learning approaches offer another route to automated diagnosis, but
classification accuracy alone is insufficient for an autonomous spacecraft.
A classifier may produce a high-confidence prediction even when the telemetry
belongs to a failure mode that was never represented in the training data.
Calibration and uncertainty estimation therefore become safety-relevant
properties rather than optional model diagnostics \cite{guo2017,hendrycks2017}.

This paper investigates an evidential diagnostic engine designed around that
requirement. Instead of treating the output of a classifier as a single
deterministic probability vector, the system assigns evidence to competing
fault hypotheses and represents the resulting uncertainty with a Dirichlet
distribution. Physical relationships within the spacecraft electrical and
thermal systems are incorporated through the feature representation and
distance metric.

The central objective is not simply to maximize classification accuracy.
The diagnostic system is expected to answer two questions simultaneously:
\begin{enumerate}
    \item Which of the known operational states best explains the observed
    telemetry?
    \item Does the observation remain sufficiently close to the known
    operational manifold for that diagnosis to be trusted?
\end{enumerate}

The second question motivates the explicit decomposition between epistemic
and aleatoric uncertainty. Epistemic uncertainty indicates that the model has
limited knowledge about the observed condition, whereas aleatoric uncertainty
reflects ambiguity or noise present in the telemetry itself.

\section{Problem Formulation and Physical Dynamics}

Let the instantaneous spacecraft telemetry vector at orbital time $t$ be
$x(t)\in\mathbb{R}^{D}$. The electrical power system is represented using
coupled electrical and thermal dynamics. The bus voltage is modeled as

\begin{equation}
V_{\mathrm{bus}}(t)
=
V_{\mathrm{oc}}(t)
-
I_{\mathrm{batt}}(t)R_0(T)
-
V_{\mathrm{RC}}(t),
\label{eq:vbus}
\end{equation}

where $V_{\mathrm{oc}}$ is the open-circuit voltage,
$I_{\mathrm{batt}}$ is battery current, $R_0(T)$ is temperature-dependent
internal resistance, and $V_{\mathrm{RC}}$ represents the transient
resistive-capacitive contribution.

The thermal state is modeled as

\begin{equation}
C_{\mathrm{th}}\frac{dT}{dt}
=
I_{\mathrm{batt}}^2R_0(T)
+
\dot{Q}_{\mathrm{exo}}
-
\epsilon\sigma S_B A
\left(T^4-T_{\mathrm{space}}^4\right),
\label{eq:thermal}
\end{equation}

where $C_{\mathrm{th}}$ is the effective thermal capacitance,
$\dot{Q}_{\mathrm{exo}}$ represents external heat input,
$\epsilon$ is emissivity, $\sigma$ is the Stefan--Boltzmann constant,
$S_BA$ denotes the effective radiating area term, and $T_{\mathrm{space}}$
is the surrounding space temperature.

The diagnostic task is to infer an operational state

\begin{equation}
y\in\mathcal{Y}=\{1,\ldots,K\},
\end{equation}

where $K=6$ and the state set contains five modeled failure modes together
with nominal operation. In addition to the class estimate, the system
produces an epistemic rejection score

\begin{equation}
u_{\mathrm{epistemic}}\in[0,1],
\end{equation}

which is intended to indicate whether the current observation lies outside
the known operational manifold $\mathcal{M}_{\mathrm{known}}\subset
\mathbb{R}^{D}$.

\section{Evidential Dirichlet Formulation}

The proposed diagnostic engine follows the evidential-learning formulation
of uncertainty-aware classification \cite{sensoy2018,malinin2018}. Instead of
assigning a fixed probability vector directly, the model represents the
categorical distribution with a Dirichlet density,

\begin{equation}
\mathrm{Dir}(\mathbf{p}\mid\boldsymbol{\alpha})
=
\frac{1}{B(\boldsymbol{\alpha})}
\prod_{k=1}^{K}p_k^{\alpha_k-1},
\label{eq:dirichlet}
\end{equation}

where $\boldsymbol{\alpha}=[\alpha_1,\ldots,\alpha_K]^T$ contains the
concentration parameters.

For each operational state, non-negative evidence is generated from the
distance between the observed telemetry and the corresponding physical
cluster. The concentration parameter is

\begin{equation}
\alpha_k=e_k+1,
\label{eq:alpha}
\end{equation}

where the evidence term is

\begin{equation}
e_k
=
\exp\left(
-\gamma
\sqrt{
(x-\mu_k)^T
\Sigma_k^{-1}
(x-\mu_k)
}
\right).
\label{eq:evidence}
\end{equation}

Here, $\mu_k$ and $\Sigma_k$ are the mean vector and covariance matrix for
operational state $k$. The parameter $\gamma$ controls how rapidly evidence
falls with increasing distance; the experiments use $\gamma=2.5$.

The expected posterior class probability is obtained from

\begin{equation}
\hat{p}_k=\frac{\alpha_k}{S},
\qquad
S=\sum_{j=1}^{K}\alpha_j.
\label{eq:posterior}
\end{equation}

This formulation allows the model to distinguish between a situation in
which one hypothesis has strong supporting evidence and a situation in which
all hypotheses have weak evidence.

\subsection{Epistemic and Aleatoric Uncertainty}

Epistemic uncertainty is associated with regions of the feature space for
which the diagnostic model has little supporting knowledge. It is calculated
from the minimum Mahalanobis distance to the known fault clusters:

\begin{equation}
u_{\mathrm{epistemic}}
=
\frac{1}
{1+\exp[-\beta(d_{\min}-\theta)]},
\label{eq:epistemic}
\end{equation}

where $d_{\min}$ is the smallest cluster distance,
$\theta=3.5\sigma$ is the empirical manifold boundary, and $\beta=1.2$
controls the transition around the boundary.

Aleatoric uncertainty is computed from the normalized Shannon entropy of the
expected posterior:

\begin{equation}
u_{\mathrm{aleatoric}}
=
\frac{
-\sum_{k=1}^{K}\hat{p}_k
\log_2(\hat{p}_k+\epsilon)
}
{\log_2(K)},
\label{eq:aleatoric}
\end{equation}

with $\epsilon$ used for numerical stability. The resulting value lies in
$[0,1]$.

This distinction is important operationally. A noisy measurement of a known
state should primarily increase aleatoric uncertainty. An observation
representing an unfamiliar physical condition should instead produce a
substantial increase in epistemic uncertainty.

\section{Experimental Methodology}

The evaluation was structured to prevent information from the test set from
affecting model selection or threshold selection. Three controls were applied.

\subsection{Scenario-Level Stratification}

Data were partitioned at the trajectory level rather than by independently
sampling individual telemetry frames. The random seeds were fixed to
\texttt{42} for training, \texttt{1337} for validation, and \texttt{2026}
for testing. This prevents adjacent samples from the same trajectory from
appearing in both training and evaluation partitions.

\subsection{Pre-Evaluation OOD Freezing}

OOD failure categories were specified before test-set scoring. The definition
of an unseen failure therefore could not be modified after observing the test
results.

\subsection{Zero Test Tuning}

The rejection threshold was selected only using validation data. The locked
threshold used throughout the reported evaluation was

\begin{equation}
\tau_{\mathrm{locked}}=0.0648.
\end{equation}

No test-set tuning was performed.

\begin{table}[t]
\centering
\caption{Telemetry Datasets: Characterization and Summary Statistics}
\label{tab:datasets}
\small
\begin{tabular}{llrrrr}
\toprule
Dataset & Channel & Unit & Mean $\pm$ Std & Min & Max\\
\midrule
NASA PCoE B0005 & Voltage & V &
$3.470\pm0.311$ & 2.542 & 4.044\\
& Current & A &
$2.000\pm0.010$ & 1.961 & 2.039\\
& Temperature & $^\circ$C &
$29.919\pm2.309$ & 23.879 & 34.822\\
& Capacity & Ah &
$1.536\pm0.187$ & 1.250 & 1.957\\
\midrule
Spacecraft EPS Sim & $V_{\mathrm{bus}}$ & V &
$32.723\pm0.465$ & 31.673 & 34.569\\
& $I_{\mathrm{batt}}$ & A &
$-10.346\pm11.817$ & -31.395 & 4.102\\
& $T_{\mathrm{core}}$ & $^\circ$C &
$11.938\pm4.349$ & 2.133 & 19.927\\
& Net Power & W &
$-342.04\pm389.56$ & -1066.9 & 132.8\\
& $R_{\mathrm{int}}$ & $\Omega$ &
$2.130\pm1.756$ & 0.001 & 13.969\\
\bottomrule
\end{tabular}
\end{table}

The NASA PCoE B0005 data contain battery-aging trajectories with capacity
decreasing from approximately $1.957$ Ah to $1.250$ Ah. The spacecraft EPS
simulation complements this dataset by representing coupled electrical and
thermal behavior over three complete LEO orbits. The simulated trajectories
include high-current discharge during eclipse periods and corresponding
thermal cycling.

\section{Experimental Results}

\subsection{Study 1: Known-Fault Classification}

The first experiment evaluated the diagnostic models on 240 held-out frames.
Table~\ref{tab:known} summarizes the results.

\begin{table}[t]
\centering
\caption{Known-Fault Benchmark on the Held-Out Test Set}
\label{tab:known}
\small
\begin{tabular}{lrrrr}
\toprule
Diagnostic Model & Accuracy & Macro-F1 & ECE & Mean Conf.\\
\midrule
Physics Rules & 0.6667 & 0.5542 & 0.2548 & 0.6708\\
Random Forest & 1.0000 & 1.0000 & 0.0308 & 0.9692\\
MLP Softmax & 1.0000 & 1.0000 & 0.0044 & 0.9956\\
Standard Mahalanobis & 1.0000 & 1.0000 & 0.0074 & 0.9926\\
Evidential Dirichlet & 0.9542 & 0.9533 & 0.0094 & 0.9561\\
\bottomrule
\end{tabular}
\end{table}

The evidential model obtains a Macro-F1 of 0.9533 and an ECE of 0.0094.
Although several deterministic baselines obtain perfect classification
accuracy on this particular held-out set, their behavior must be considered
alongside calibration and OOD performance. Relative to the physics-rule
baseline, the proposed method shows a statistically significant improvement
with $p=3.94\times10^{-16}$ and Cohen's $d=+0.634$.

\subsection{Study 2: Uncertainty Disentanglement}

The second study examined whether the two uncertainty measures responded to
different sources of variation. Across 1,400 evaluated frames, increasing
sensor noise was positively associated with aleatoric uncertainty
($\rho_{\mathrm{Spearman}}=0.2955$, $p=1.37\times10^{-17}$).

At a severe telemetry noise level of $\sigma=0.20$, epistemic uncertainty
remained approximately 0.092, substantially below the rejection reference
level of 0.45. Thus, noise in an otherwise familiar operating condition did
not automatically cause the system to classify the observation as unknown.

For novel, unmodeled faults, epistemic uncertainty increased to 1.000,
corresponding to a reported separation ratio of approximately 10.94.

\subsection{Study 3: Out-of-Distribution Detection}

The third experiment compared the proposed epistemic gating mechanism with
several OOD baselines. The benchmark contained 600 in-distribution frames and
600 OOD frames.

\begin{table}[t]
\centering
\caption{Out-of-Distribution Detection Benchmark}
\label{tab:ood}
\small
\begin{tabular}{lrrr}
\toprule
Method & AUROC & AUPRC & FPR @ 95\% TPR\\
\midrule
MLP Softmax Inverted & 0.4313 & 0.6135 & 1.0000\\
Random Forest Variance & 0.9384 & 0.9426 & 0.3200\\
Evidential Dirichlet & 0.9422 & 0.9516 & 0.3583\\
Isolation Forest & 0.9815 & 0.9823 & 0.1150\\
\bottomrule
\end{tabular}
\end{table}

At the locked validation threshold $\tau_{\mathrm{locked}}=0.0648$, the
evidential gating mechanism achieved a TPR of 76.67\% at an FPR of 4.50\%.
The inverted MSP baseline failed to separate the two regimes effectively,
with an AUROC of only 0.4313. This result illustrates the danger of relying
on softmax confidence alone when the input may fall outside the training
distribution.

The Isolation Forest baseline achieved the highest OOD AUROC in this
benchmark. Consequently, the results do not support the claim that
evidential inference is universally superior for OOD detection. Instead, they
show that the proposed approach provides a physically grounded uncertainty
signal while remaining competitive with the evaluated alternatives.

\subsection{Study 4: Telemetry Noise Robustness}

Telemetry noise was increased progressively from $\sigma=0.00$ to
$\sigma=0.25$. Across this range, the evidential engine maintained a
statistically significant advantage over the physics-rule baseline, with
$p=5.77\times10^{-6}$ and a mean Macro-F1 advantage of $+0.3152$.

For the flight-representative noise range of $\sigma\leq0.05$, the false-OOD
alarm rate remained below 7.7\%. This behavior is consistent with the intended
uncertainty decomposition: moderate observation noise increases aleatoric
uncertainty without forcing a large epistemic rejection signal.

\subsection{Study 5: Systematic Component Ablation}

Ablation experiments were used to determine which components were responsible
for the observed behavior.

\begin{table}[t]
\centering
\caption{Quantitative Component Ablation}
\label{tab:ablation}
\small
\begin{tabular}{lrrrr}
\toprule
Configuration & Macro-F1 & ECE & OOD AUROC & $R_{\mathrm{int}}$ F1\\
\midrule
M0 Full & 0.9533 & 0.0179 & 0.8692 & 1.0000\\
M1 No Evidential & 0.9533 & 0.0179 & 0.7895 & 1.0000\\
M2 Euclidean Metric & 0.7132 & 0.2202 & 0.6870 & 0.0000\\
M3 No Physics Priors & 0.2047 & 0.0465 & 0.7811 & 0.3493\\
M4 Static Features & 0.9491 & 0.1317 & 0.9070 & 0.9873\\
\bottomrule
\end{tabular}
\end{table}

The Euclidean-metric ablation (M2) produced the clearest degradation in
internal-resistance-spike detection, reducing its F1 score to zero. This
indicates that the covariance-aware Mahalanobis geometry is important for
capturing relationships among the telemetry channels.

Removing the evidential component (M1) reduced OOD AUROC from 0.8692 to
0.7895 in the ablation configuration, a decrease of approximately 8
percentage points. Removing the physics priors (M3) caused a much larger
drop in known-fault Macro-F1, to 0.2047. Static features (M4) preserved most
of the classification performance but substantially worsened calibration.

\section{Identified Failure Modes and Limitations}

The experiments also reveal boundaries that should be considered before using
the method in a flight-critical setting.

\subsection{Compound-Fault Centroid Cancellation}

When two failures occur simultaneously, their effects on different telemetry
channels can point in opposing directions. In the tested solar-loss plus
thermal-surge case, the resulting feature vector can fall into an
intermediate region between modeled fault clusters. The measured epistemic
uncertainty in this regime was $0.301\pm0.092$. Consequently, the metric
representation can underestimate the unfamiliarity of some compound faults.

\subsection{Magnitude-Invariance Blind Spot}

The use of magnitude-based current features can remove the sign information
needed to recognize a polarity reversal. In the sensor sign-inversion
experiment, the upstream transformation based on $|I_{\mathrm{batt}}|$
masked the reversal and reduced the catch rate to 6.7\%.

These limitations suggest two direct directions for future work: explicitly
preserving signed telemetry features and introducing compound-fault training
or structured multi-fault representations.

\section{Conclusion}

This paper evaluated an evidential Dirichlet approach for uncertainty-aware
fault diagnosis in autonomous spacecraft health management. The method
combines physically motivated telemetry representations with a
covariance-aware distance measure and an evidential output layer. The
resulting model provides both a diagnosis over known operational modes and
an uncertainty signal intended to identify observations that do not belong
to the known operating manifold.

On the reported held-out benchmark, the system achieved a Macro-F1 of 0.9533
and an ECE of 0.0094. Telemetry noise primarily increased aleatoric
uncertainty, while novel failure conditions produced a large epistemic
response. The resulting OOD detector reached an AUROC of 0.9422 and an AUPRC
of 0.9516, substantially outperforming the evaluated MSP baseline.

At the same time, the experiments show that uncertainty estimation does not
remove the need for careful feature engineering and physical modeling.
Compound faults can produce metric cancellation, while magnitude-only
representations can hide polarity changes. These failure boundaries are
therefore part of the contribution of the study: an uncertainty-aware
diagnostic architecture should be evaluated not only by its successful cases,
but also by the physical situations in which its assumptions break down.

Future versions of AstraHeal should preserve signed sensor information,
support explicit compound-fault hypotheses, and evaluate the uncertainty
gating mechanism on broader spacecraft telemetry and independently generated
fault scenarios.

\begin{thebibliography}{00}

\bibitem{chien2005}
S. Chien, R. Sherwood, D. Tran, B. Cichy, D. Zhu, R. Castano, A. Davies,
D. Mandl, S. Frye, B. Trout, et al.,
``Autonomous sciencecraft experiment on the EO-1 mission,''
\emph{Journal of Aerospace Computing, Information, and Communication},
vol. 2, no. 4, pp. 196--228, 2005.

\bibitem{muscettola1998}
N. Muscettola, P. P. Nayak, B. Pell, and B. C. Williams,
``Remote agent: To boldly go where no AI has gone before,''
\emph{Artificial Intelligence}, vol. 103, nos. 1--2, pp. 5--47, 1998.

\bibitem{williams1996}
B. C. Williams and P. P. Nayak,
``A model-based approach to reactive self-configuring systems,''
in \emph{Proc. AAAI Conference on Artificial Intelligence},
vol. 2, 1996, pp. 971--978.

\bibitem{guo2017}
C. Guo, G. Pleiss, Y. Sun, and K. Q. Weinberger,
``On calibration of modern neural networks,''
in \emph{Proc. International Conference on Machine Learning (ICML)},
2017, pp. 1321--1330.

\bibitem{hendrycks2017}
D. Hendrycks and K. Gimpel,
``A baseline for detecting misclassified and out-of-distribution examples
in neural networks,''
in \emph{Proc. International Conference on Learning Representations (ICLR)},
2017.

\bibitem{sensoy2018}
M. Sensoy, L. Kaplan, and M. Kandemir,
``Evidential deep learning to quantify classification uncertainty,''
in \emph{Advances in Neural Information Processing Systems (NeurIPS)},
vol. 31, pp. 3179--3189, 2018.

\bibitem{malinin2018}
A. Malinin and M. Gales,
``Predictive uncertainty estimation via prior networks,''
in \emph{Advances in Neural Information Processing Systems (NeurIPS)},
vol. 31, pp. 7047--7058, 2018.

\end{thebibliography}

\end{document}
